\introduction

I once attended a talk by the poet Craig Raine on his libretto
to the opera \textit{The Electrification of the Soviet Union}
(music by Nigel Osbourne, stage direction by Peter Sellars,
libretto published by Faber and Faber in 1986).
Raine had been approached by Osbourne to write a libretto
based on Pasternak's novella \textit{The Last Summer}, this being
chosen by Osbourne because it was (in Osbourne's words) `The
most musical book he knew.'  Raine was completely taken off-guard
by this.  He did not at the time, nor later when he was giving
the talk on his completed libretto, know what Osbourne had meant.
To me, the musicality of \textit{The Last Summer} shines through
and is one of the main delights of this work, yet an excellent poet
such as Raine, who surely has some well-developed concept of music, did not see it
at all.

This is not an isolated occurrence. T.S.~Eliot wrote a well-known
essay entitled \textit{The Music of Poetry}\footnote{In: On Poetry and
Poets, Faber and Faber 1957.}, in which he never explains what he
means by `music' in poetry.  The closest he gets to a definition
seems to be the highly circular,
\begin{quote}
\ldots a `musical poem' is a poem which has a musical pattern
of sound and a musical pattern of the secondary meanings of the
words which compose it, and that these  two patterns are indissoluble and one.
\end{quote}

Thus we learn of Eliot's opinion (which I wholeheartedly agree with)
that the `music' of poetry is more than the rhymes, assonances or alliteration
that may be there and more than the rhythms of feet and metre which we hear.
Meaning plays an essential role too.  But nowhere in this essay
does Eliot explain what he means by `music,' except by appealing to
the reader's preconceived idea of what is meant by music \textit{qua} music.

We get a bit closer to what Osbourne was driving at a little later
when Eliot writes that `a play of Shakespeare is a very complex musical
structure.'  In other words, the musicality of a piece of literature
of whatever length is an attribute that incorporates its entirity and overall structure.
These are further hints that the music of poetry or literature is not about
rhyme or alliteration or metre but something more, perhaps something that
introduces, develops and concludes a theme in the way musical form does.
Music that is music must attempt
to be more than a catching ditty or mere tuneful phrase.  In this respect,
Pasternak's novella \textit{The Last Summer} is highly musical, at least in
my opinion as well as Osbourne's.

In Eliot's very last paragraph, where perhaps a summing-up is expected,
he introduces new ideas for his essay that hint at possible future developments of
his thesis that align closely to what has been said about musicality involving
the entirity of a piece. 
\begin{quote}
I think that a poet may gain much from the study of music:
how much technical knowledge of musical form is desirable I do not know, for I have
not that technical knowledge myself.
\end{quote}
and
\begin{quote}
The use of recurrent themes is as natural to poetry as to music.  There are possibilities
for verse which bear some analogy to the development of a theme by different groups of instruments;
there are possibilities of transitions in a poem comparable to the different movements
of a symphony or a quartet; there are possibilities of contrapuntal arrangement of subject-matter.
[\ldots] More than this I cannot say, but must leave the matter here to those who have had a musical education.
\end{quote}

I have also, for some time, thought that a closer connection between literature and
music would be beneficial to both.  This volume is a modest contribution in this direction.

The \textit{Well Tempered Clavier} by J.S.Bach, usually referred to informally as `The 48,'
is a set of two books, each book containing a prelude and fugue in each of the 24 possible
major and minor keys.  It was written to showcase the tuning that had been recently
introduced on keyboard instruments such as organs and harpsichords, in
which every octave is divided into 12 equal intervals (semitones) and therefore
as a result the keyboard can play equally well in every possible key. This is of course
nowadays the usual tuning on a piano.

Bach's 48~successfully achieves all of this and the idea has frequently been copied, for example
by Chopin in his \textit{Preludes}. But it achieves
much more than this.  It is a showcase for composition: each key inviting the composer
to write in a different style, always within the constraints of an introductory
prelude followed by a contrapuntal fugue, the word literally meaning `flight'.  Furthermore,
composers, performers and those listening to music have learnt to associate particular moods
or emotions with each key.  For example, the key of E\fl\ major is associated with
the `heroic,' as used for example in Beethoven's \textit{Eroica} symphony or
Strauss's tone poem \textit{Ein Heldenleben}.  Whether these moods associated with keys
is an innate part of our psychological make-up or a synthetic product of the
way composers write music is a moot point, but there is no doubt that pieces of music
in different keys sound very different in terms of mood and emotion.

The same applies to literature.  If anything can convey emotion and mood, literature surely can,
and there are considerably more traits available to the author or poet than the 24~keys
available to Bach.

With all this in mind, this volume is a small set of preludes and fugues in literature. Each author has been 
asked to focus on one or more moods or emotions and a person, object or place that conveys those traits, and 
write an introductory `prelude' setting out the main theme, person, object or place, followed by a `fugue' or 
story or poem involving that.  Of course, to what extent these pieces are musical is up to the reader to judge, 
but we have all found the exercise useful in expanding and developing style and technique in writing and we hope 
you will enjoy them all.
